\documentclass[11pt,oneside]{article}

\usepackage[T1]{fontenc}
\usepackage{lmodern}
\usepackage{microtype}
\usepackage[a4paper,margin=28mm,headheight=15pt]{geometry}
\usepackage{amsmath,amssymb,amsthm,mathtools}
\usepackage{booktabs,longtable,array}
\usepackage{enumitem,multicol}
\usepackage{xcolor}
\usepackage{hyperref}
\usepackage[nameinlink,noabbrev]{cleveref}
\usepackage{fancyhdr}
\usepackage{listings}
\usepackage{url}

\definecolor{deepblue}{RGB}{22,55,92}
\definecolor{softblue}{RGB}{235,242,249}
\definecolor{darkgreen}{RGB}{20,95,65}
\hypersetup{
  colorlinks=true,
  linkcolor=deepblue,
  citecolor=darkgreen,
  urlcolor=deepblue,
  pdftitle={Integral and Transform Frontiers for the Fabius and Rvachev Functions},
  pdfauthor={OpenAI GPT-5.6 Pro}
}

\pagestyle{fancy}
\fancyhf{}
\fancyhead[L]{Fabius--Rvachev Integral Frontiers}
\fancyhead[R]{\thepage}
\renewcommand{\headrulewidth}{0.4pt}

\setlength{\parindent}{0pt}
\setlength{\parskip}{0.58em}
\setlist[itemize]{topsep=0.25em,itemsep=0.18em,leftmargin=1.6em}
\setlist[enumerate]{topsep=0.25em,itemsep=0.18em,leftmargin=1.8em}
\allowdisplaybreaks

\newtheorem{theorem}{Theorem}[section]
\newtheorem{proposition}[theorem]{Proposition}
\newtheorem{lemma}[theorem]{Lemma}
\newtheorem{corollary}[theorem]{Corollary}
\newtheorem{conjecture}[theorem]{Conjecture}
\theoremstyle{definition}
\newtheorem{definition}[theorem]{Definition}
\newtheorem{example}[theorem]{Example}
\theoremstyle{remark}
\newtheorem{remark}[theorem]{Remark}

\newcommand{\R}{\mathbb R}
\newcommand{\C}{\mathbb C}
\newcommand{\N}{\mathbb N}
\newcommand{\E}{\mathbb E}
\newcommand{\Prob}{\mathbb P}
\newcommand{\ii}{\mathrm i}
\newcommand{\dd}{\,\mathrm d}
\newcommand{\upf}{\operatorname{up}}
\newcommand{\sinc}{\operatorname{sinc}}
\newcommand{\falling}[2]{#1^{\underline{#2}}}
\newcommand{\rising}[2]{#1^{\overline{#2}}}
\newcommand{\repo}{\texttt{ProveIt}}
\newcommand{\newresult}{\par\medskip\noindent\textcolor{deepblue}{\textbf{Frontier contribution.}}\ }

\lstset{
  basicstyle=\ttfamily\small,
  columns=fullflexible,
  frame=single,
  breaklines=true,
  showstringspaces=false
}

\title{\bfseries Integral and Transform Frontiers\\[0.25em]
for the Fabius, Inverse Fabius, and Rvachev Functions}
\author{A repository-aware analytic and computational investigation}
\date{27 August 2026}

\begin{document}
\maketitle

\begin{abstract}
This report develops a systematic integral calculus for the Fabius function $F$, Rvachev's compactly supported function $\upf$, and the inverse Fabius function $G=F^{-1}$.  It begins with a recursive audit of the current TeX corpus in the \repo{} repository, including the main monograph, the consolidated non-formalized frontier document, imported papers, current draft packages, and historical archives.  Existing derivative, moment, Fourier-product, Lambert--$W$, Bell-polynomial, Bernoulli-polynomial, and inverse-endpoint results are treated as baselines rather than reannounced as new.

The principal new results are: a finite dyadic closed form for $\int x^pF(x)\dd x$ when $p\in\N_0$; an absolutely and locally uniformly convergent dyadic series for arbitrary complex powers, including negative and fractional exponents; two-parameter recurrences for fractional Volterra integrals; complementary Mellin transforms controlled by an entire complex moment function; an entire Newton interpolation of complex moments from rational integer moments; exact local primitives and all-order derivative formulas for the inverse Fabius function; a sharp $L^q$ threshold for $G'$; refined inverse-primitive endpoint asymptotics carrying the repository's log-periodic Lambert phase; Stieltjes, logarithmic-potential, autocorrelation, and order-statistic transforms; and a rational Jacobi continued fraction generated by the Fabius moments.  Numerical experiments based on finite Thue--Morse splines independently check the main formulas.  Unproved conjectures and a Lean formalization roadmap are included.
\end{abstract}

\tableofcontents
\clearpage

\section{Scope, corpus audit, and novelty protocol}

\subsection{Repository review}

The live directory
\begin{center}
\url{https://github.com/VladimirReshetnikov/ProveIt/tree/main/Analysis/FabiusFunction/docs}
\end{center}
was recursively audited on 27 August 2026.  All discoverable TeX paths were inventoried.  Current synthesis documents and research drafts were examined through their raw TeX, imported papers were checked for the definitions and results relevant to integration, and historical copies were compared through their content indexes so that superseded duplicates were not counted as independent evidence.

The main baselines were:
\begin{itemize}
\item \texttt{Fabius\_Function\_and\_Rvachev\_Up.tex}, the principal monograph;
\item \texttt{non-formalized-research-frontiers.tex}, the consolidated frontier synthesis;
\item the Lean walkthrough, glossary, and non-elementarity report;
\item the imported TeX sources of Arias de Reyna, Haugland, Aistleitner--Hofer--Larcher, T\'oth, and related papers;
\item the archive index and its historical report families.
\end{itemize}

The current draft inbox contained the following additional packages:
\begin{multicols}{2}\small
\begin{itemize}
\item \path{Fabius_Arithmetic_Rays_Frontier_Report}
\item \path{Fabius_Dyadic_Self_Sampling_Frontier_Package}
\item \path{Fabius_Half_Integer_Spectral_Frontier_Report}
\item \path{Fabius_Inverse_Frontier_Report_Source_and_PDF}
\item \path{Fabius_Newton_Rvachev_Frontier_Report}
\item \path{Fabius_Rvachev_Frontier_Report-2}
\item \path{Fabius_Rvachev_Frontier_Report}
\item \path{Fabius_Rvachev_New_Frontiers}
\item \path{Fabius_Rvachev_Thue_Morse_Frontier_Results}
\item \path{Spectral_Arithmetic_Pascal_Rvachev_Hierarchy}
\item \path{Thue_Morse_Formula_Atlas}
\item \path{fabius_frontier_dyadic_inverse_barnes_report}
\item \path{fabius_frontier_new_results}
\item \path{fabius_frontier_results}
\item \path{fabius_frontier_results_bundle}
\item \path{fabius_frontier_spectral_endpoint_report_bundle}
\item \path{rvachev_up_fourier_decay}
\end{itemize}
\end{multicols}

The primitive identity $F(x/2)=\int_0^xF(t)\dd t$ is implicit in the audited spectral hierarchy and therefore is used here as a baseline.  Likewise, the repository already contains extensive work on integer moments, dyadic values, Fourier sinc products, endpoint asymptotics, inverse asymptotics, and Bell-polynomial reversion.  The novelty claims below are deliberately narrower: ``new'' means not located explicitly in this audited corpus or in the targeted external literature search, not an absolute historical-priority claim.

\subsection{Main questions addressed}

The report studies:
\begin{itemize}
\item normalized and indefinite antiderivatives of $F$, $\upf$, $1-F$, $F(1-x)$, and $G$;
\item polynomial, negative, fractional, complex, exponential, logarithmic, and resolvent weights;
\item integer and fractional repeated integration, and all-order differentiation of the inverse;
\item Mellin, Laplace, Fourier, Stieltjes/Cauchy, logarithmic-potential, and correlation transforms;
\item products and compositions such as $F^a(1-F)^b$, $\upf(2G-1)^r$, and shifted products of $\upf$;
\item exact arithmetic consequences, numerical verification, conjectures, and formalization targets.
\end{itemize}

\section{Definitions and probabilistic normalization}

\subsection{The bounded Fabius function and Rvachev's function}

Let $F:\R\to[0,1]$ be the bounded Fabius distribution function, with
\begin{equation}\label{eq:F-boundary}
 F(x)=0\quad(x\le0),
 \qquad F(x)=1\quad(x\ge1),
\end{equation}
and, on $0\le x\le1/2$,
\begin{equation}\label{eq:F-differential}
 F'(x)=2F(2x).
\end{equation}
It obeys
\begin{equation}\label{eq:F-symmetry}
 F(x)+F(1-x)=1,
 \qquad 0\le x\le1.
\end{equation}
Rvachev's function is
\begin{equation}\label{eq:up-definition}
 \upf(u)=
 \begin{cases}
 F(1-|u|),&|u|\le1,\\
 0,&|u|>1.
 \end{cases}
\end{equation}
It is even, nonnegative, $C^\infty$, supported on $[-1,1]$, and normalized by
$\int_{-1}^{1}\upf(u)\dd u=1$.  The two functions satisfy
\begin{equation}\label{eq:Fprime-up}
 F'(x)=2\upf(2x-1),
 \qquad 0<x<1,
\end{equation}
and, on $[0,1]$,
\begin{equation}\label{eq:up-survival}
 \upf(x)=1-F(x)=F(1-x).
\end{equation}

The repository also uses a signed global extension $\mathcal F$ that agrees with $F$ on $[0,1]$ and satisfies
\begin{equation}\label{eq:signed-differential}
 \mathcal F'(x)=2\mathcal F(2x),
 \qquad x\in\R.
\end{equation}
It is indispensable when derivatives move the argument outside $[0,1]$.

\subsection{Random-series model}

Let $U_1,U_2,\ldots$ be independent uniform random variables on $[0,1]$ and set
\begin{equation}\label{eq:X-random-series}
 X=\sum_{j=1}^{\infty}\frac{U_j}{2^j}.
\end{equation}
Then $F$ is the distribution function of $X$.  If $V_j=2U_j-1$ and
\begin{equation}\label{eq:Y-random-series}
 Y=2X-1=\sum_{j=1}^{\infty}\frac{V_j}{2^j},
\end{equation}
then $Y$ has density $\upf$.  In particular,
\begin{equation}\label{eq:up-tail-probability}
 \upf(t)=\Prob(X>t),
 \qquad 0\le t\le1.
\end{equation}
This survival-function interpretation is the quickest route to many definite integrals.

Let
\begin{equation}\label{eq:D-definition}
 D(z)=\E X^z,
 \qquad d_n=D(n).
\end{equation}
Because the density is infinitely flat at $0$, all negative moments exist and $D$ extends to an entire function; this is proved in \cref{prop:D-entire}.

\subsection{Endpoint flatness and the inverse}

Every derivative of $F$ vanishes at $0$, every derivative of $1-F$ vanishes at $1$, and every derivative of $\upf-1$ vanishes at $0$.  Thus, for each $M\in\N$,
\begin{equation}\label{eq:flatness-bounds}
 F(x)=O(x^M),\quad F'(x)=O(x^M)
 \qquad(x\downarrow0).
\end{equation}
These bounds justify complex powers, repeated integration by parts, and normal convergence arguments that would fail for an ordinary density with algebraic endpoint behavior.

The restriction of $F$ to $[0,1]$ is strictly increasing.  Its inverse
\begin{equation}\label{eq:G-definition}
 G:(0,1)\to(0,1),
 \qquad G=F^{-1},
\end{equation}
obeys
\begin{equation}\label{eq:G-symmetry}
 G(1-y)=1-G(y).
\end{equation}

\section{Three master integration principles}

\begin{theorem}[CDF--survival integration]\label{thm:master-H}
Let $H$ be absolutely continuous on $[0,1]$ and suppose the following integrals exist.  Then
\begin{align}
 \int_0^1H'(t)\upf(t)\dd t
 &=\E\bigl[H(X)-H(0)\bigr],\label{eq:master-up}\\
 \int_0^1H'(t)F(t)\dd t
 &=H(1)-\E H(X).\label{eq:master-F}
\end{align}
\end{theorem}

\begin{proof}
For \eqref{eq:master-up}, use \eqref{eq:up-tail-probability} and Fubini:
\[
 \int_0^1H'(t)\Prob(X>t)\dd t
 =\E\int_0^XH'(t)\dd t.
\]
Equation \eqref{eq:master-F} follows either similarly or by adding the two formulas and using $F+\upf=1$.
\end{proof}

\begin{theorem}[Quantile push-forward]\label{thm:quantile-pushforward}
For every Borel function $\Phi$ for which either side is absolutely integrable,
\begin{equation}\label{eq:quantile-pushforward}
 \int_0^1\Phi(G(y))\dd y
 =\int_0^1\Phi(x)\dd F(x)
 =\E\Phi(X).
\end{equation}
\end{theorem}

\begin{proof}
This is the probability-integral transform, or directly the substitution $y=F(x)$.
\end{proof}

\begin{theorem}[Universal beta family]\label{thm:beta-family}
For $\Re a>-1$ and $\Re b>-1$,
\begin{equation}\label{eq:beta-family}
 \int_0^1F(x)^a\bigl(1-F(x)\bigr)^bF'(x)\dd x
 =\mathrm B(a+1,b+1).
\end{equation}
Equivalently,
\begin{equation}\label{eq:beta-up-family}
 \int_0^1F(x)^a\upf(x)^b\upf(2x-1)\dd x
 =\frac12\mathrm B(a+1,b+1).
\end{equation}
\end{theorem}

\begin{proof}
Put $y=F(x)$ and use \eqref{eq:Fprime-up}.
\end{proof}

The point of \cref{thm:master-H,thm:quantile-pushforward} is not merely probabilistic elegance.  They convert any tractable ordinary antiderivative $H$ into an exact Fabius or Rvachev integral.  Monomials give moments, exponentials give Laplace products, logarithms give entropy and logarithmic potentials, and resolvents give Stieltjes transforms.

\section{Integer-order primitives and derivative hierarchies}

\subsection{Normalized primitives of $F$}

From \eqref{eq:F-differential},
\begin{equation}\label{eq:first-primitive-F}
 \frac{\dd}{\dd x}F(x/2)=F(x),
 \qquad 0\le x\le1.
\end{equation}
Thus $F(x/2)$ is the primitive normalized to vanish at $x=0$.

For $n\in\N$, write
\begin{equation}\label{eq:In-definition}
 I_0^nf(x)=\frac1{(n-1)!}\int_0^x(x-t)^{n-1}f(t)\dd t.
\end{equation}

\begin{theorem}[Exact $n$th primitive of $F$]\label{thm:nth-primitive-F}
For $n\in\N$ and $0\le x\le1$,
\begin{equation}\label{eq:nth-primitive-F}
\boxed{
 I_0^nF(x)=2^{\binom n2}F\!\left(\frac{x}{2^n}\right).}
\end{equation}
Every other $n$th antiderivative differs from the right side by a polynomial of degree at most $n-1$.
\end{theorem}

\begin{proof}
The case $n=1$ is \eqref{eq:first-primitive-F}.  If the formula holds at order $n$, then
\[
 I_0^{n+1}F(x)=2^{\binom n2}\int_0^xF(t/2^n)\dd t
 =2^{\binom n2+n}F(x/2^{n+1}),
\]
which is the required triangular exponent.
\end{proof}

For the bounded extension and $x\ge1$, the $n$th derivative is $1$, so matching derivatives at $x=1$ gives
\begin{equation}\label{eq:nth-primitive-F-tail}
 I_0^nF(x)=\frac{(x-1)^n}{n!}
 +\sum_{k=0}^{n-1}
 \frac{2^{\binom{n-k}{2}}F(2^{-(n-k)})}{k!}(x-1)^k.
\end{equation}
For the signed global extension, a single formula is valid on all of $\R$:
\begin{equation}\label{eq:signed-nth-antiderivative}
 \frac{\dd^n}{\dd x^n}
 \left[2^{\binom n2}\mathcal F(x/2^n)\right]
 =\mathcal F(x).
\end{equation}

\subsection{Derivatives and primitives of $\upf$}

The global CDF of the centered density is
\begin{equation}\label{eq:up-first-primitive}
 \int_{-\infty}^{x}\upf(t)\dd t
 =F\!\left(\frac{x+1}{2}\right).
\end{equation}
Consequently:

\begin{theorem}[Left $n$th primitive of $\upf$]\label{thm:nth-primitive-up}
For $n\in\N$ and $x\le1$,
\begin{equation}\label{eq:nth-primitive-up}
\boxed{
 \frac1{(n-1)!}\int_{-\infty}^{x}(x-t)^{n-1}\upf(t)\dd t
 =2^{\binom n2}F\!\left(\frac{x+1}{2^n}\right).}
\end{equation}
For $x\ge1$, it becomes the moment polynomial
\begin{equation}\label{eq:nth-primitive-up-tail}
 \sum_{k=0}^{n-1}
 \frac{(-1)^k\mu_k}{k!(n-1-k)!}x^{n-1-k},
 \qquad
 \mu_k=\int_{-1}^{1}t^k\upf(t)\dd t.
\end{equation}
\end{theorem}

\begin{proof}
Equation \eqref{eq:up-first-primitive} is the case $n=1$.  Induct as in \cref{thm:nth-primitive-F}, now mapping the lower endpoint $-1$ to $0$.  For $x\ge1$, expand $(x-t)^{n-1}$ over the compact support.
\end{proof}

On $0\le x\le1$ one also has the convenient centered primitive
\begin{equation}\label{eq:local-up-primitive}
 \int_0^x\upf(t)\dd t
 =x-F(x/2)
 =F\!\left(\frac{x+1}{2}\right)-\frac12.
\end{equation}

The forward derivative hierarchy is
\begin{equation}\label{eq:forward-derivative-hierarchy}
 \mathcal F^{(n)}(x)=2^{\binom{n+1}{2}}\mathcal F(2^nx),
\end{equation}
and, for $-1\le x\le1$,
\begin{equation}\label{eq:up-derivative-hierarchy}
 \upf^{(n)}(x)=2^{\binom{n+1}{2}}
 \mathcal F\bigl(2^n(x+1)\bigr).
\end{equation}
The antiderivative hierarchy \eqref{eq:nth-primitive-F} is therefore the exact inverse triangular scaling of the derivative hierarchy.

\section{Dyadic antiderivatives with polynomial, fractional, and complex weights}

This section contains the main antiderivative results of the report.  The elementary dilation primitive in \cref{eq:first-primitive-F} becomes substantially more useful after one allows a monomial weight and then analytically continues the exponent.

\subsection{A finite closed form for nonnegative integer powers}

For $0<x\le1$ and $\nu\in\C$ put
\begin{equation}\label{eq:A-nu-def}
 A_\nu(x)=\int_0^x t^\nu F(t)\dd t,
 \qquad t^\nu=e^{\nu\log t}.
\end{equation}
Endpoint flatness makes this integral locally uniformly convergent for every $\nu\in\C$; hence $A_\nu(x)$ is entire as a function of $\nu$.

\begin{lemma}[One-step dyadic reduction]\label{lem:A-recurrence}
For $0<x\le1$ and every $\nu\in\C$,
\begin{equation}\label{eq:A-recurrence}
 A_\nu(x)=x^\nu F(x/2)-\nu 2^\nu A_{\nu-1}(x/2).
\end{equation}
\end{lemma}

\begin{proof}
On $[0,1]$, $\frac{\dd}{\dd t}F(t/2)=F(t)$.  Integration by parts gives
\[
 A_\nu(x)=\bigl[t^\nu F(t/2)\bigr]_{0}^{x}
 -\nu\int_0^x t^{\nu-1}F(t/2)\dd t.
\]
The lower boundary vanishes for all complex $\nu$ because $F$ is flat at $0$.  In the remaining integral put $t=2u$.
\end{proof}

\begin{theorem}[Closed monomial antiderivative]\label{thm:integer-weight-antiderivative}
For every $p\in\N_0$ and $0\le x\le1$,
\begin{equation}\label{eq:integer-weight-antiderivative}
\boxed{
 \int_0^x t^pF(t)\dd t
 =\sum_{j=0}^{p}(-1)^j\frac{p!}{(p-j)!}
 2^{\binom{j+1}{2}}x^{p-j}
 F\!\left(\frac{x}{2^{j+1}}\right).}
\end{equation}
In particular,
\begin{align}
\int_0^x tF(t)\dd t
 &=xF(x/2)-2F(x/4),\label{eq:A1}\\
\int_0^x t^2F(t)\dd t
 &=x^2F(x/2)-4xF(x/4)+16F(x/8),\label{eq:A2}\\
\int_0^x t^3F(t)\dd t
 &=x^3F(x/2)-6x^2F(x/4)+48xF(x/8)-384F(x/16).
\end{align}
For $x\ge1$, the bounded extension is obtained by adding the elementary tail
\begin{equation}\label{eq:A-tail}
 A_\nu(x)=A_\nu(1)+
 \begin{cases}
 \dfrac{x^{\nu+1}-1}{\nu+1},&\nu\ne-1,\\[0.7em]
 \log x,&\nu=-1,
 \end{cases}
\end{equation}
where the first quotient has a removable singularity at $\nu=-1$.
\end{theorem}

\begin{proof}
Iterate \eqref{eq:A-recurrence}.  At the $j$th step the falling factorial contributes $p!/(p-j)!$.  The powers of $2$ telescope to
\[
 2^{jp-j(j-1)/2}\,2^{-j(p-j)}=2^{j(j+1)/2}.
\]
The process terminates because $\falling{p}{p+1}=0$.  Formula \eqref{eq:A-tail} follows from $F=1$ on $[1,\infty)$.
\end{proof}

\begin{remark}
At $x=1$, integration by parts against the probability measure $\dd F$ gives
\begin{equation}\label{eq:integer-definite-moment}
 \int_0^1t^pF(t)\dd t=\frac{1-d_{p+1}}{p+1},
 \qquad
 \int_0^1t^p\upf(t)\dd t=\frac{d_{p+1}}{p+1}.
\end{equation}
Since the integer moments $d_n$ are rational, all these definite integrals are rational.
\end{remark}

\subsection{A convergent dyadic series for arbitrary complex powers}

The recurrence does not terminate for a nonintegral exponent, but endpoint flatness converts it into a genuinely convergent series rather than merely a formal asymptotic expansion.

\begin{theorem}[Complex-weight dyadic antiderivative series]\label{thm:complex-weight-series}
Fix $0<x\le1$.  For every $\nu\in\C$,
\begin{equation}\label{eq:complex-weight-series}
\boxed{
 A_\nu(x)=
 \sum_{j=0}^{\infty}(-1)^j\falling{\nu}{j}
 2^{\binom{j+1}{2}}x^{\nu-j}
 F\!\left(\frac{x}{2^{j+1}}\right).}
\end{equation}
The series converges absolutely and locally uniformly in $\nu$.  After $N$ terms the exact remainder is
\begin{equation}\label{eq:complex-weight-remainder}
 R_N=(-1)^N\falling{\nu}{N}
 2^{N\nu-\binom N2}A_{\nu-N}(x/2^N).
\end{equation}
Consequently the sum is entire in $\nu$ and may be differentiated termwise on compact subsets of the $\nu$-plane.
\end{theorem}

\begin{proof}
Iterating \eqref{eq:A-recurrence} $N$ times gives the displayed partial sum and \eqref{eq:complex-weight-remainder}.  It remains to prove $R_N\to0$.  Flatness means that for every $M\in\N$ there is $C_M$ such that $0\le F(v)\le C_Mv^M$ on $[0,1]$.  From \cref{thm:nth-primitive-F},
\begin{align}
 2^{\binom N2}F(u/2^N)
 &=\frac1{(N-1)!}\int_0^u(u-v)^{N-1}F(v)\dd v\notag\\
 &\le C_M\frac{M!}{(N+M)!}u^{N+M}.\label{eq:flatness-factorial-bound}
\end{align}
After changing variables in \eqref{eq:complex-weight-remainder},
\[
 R_N=(-1)^N\falling{\nu}{N}
 \int_0^x u^{\nu-N}
 \bigl[2^{\binom N2}F(u/2^N)\bigr]\dd u.
\]
For $\nu$ in a fixed compact set, choose $M$ so that $\Re\nu+M>-1$ uniformly.  Bound \eqref{eq:flatness-factorial-bound} then gives
\[
 |R_N|\le C_{M,K,x}\frac{|\falling{\nu}{N}|}{(N+M)!}.
\]
The quotient tends to zero faster than a chosen inverse power of $N$, locally uniformly in $\nu$; at nonnegative integral $\nu$ it eventually vanishes identically.  This also proves normal convergence.
\end{proof}

\begin{corollary}[Negative powers]\label{cor:negative-powers}
For $m\in\N$ and $0<x\le1$,
\begin{equation}\label{eq:negative-power-series}
 \int_0^x\frac{F(t)}{t^m}\dd t
 =\sum_{j=0}^{\infty}\rising{m}{j}
 2^{\binom{j+1}{2}}x^{-m-j}
 F\!\left(\frac{x}{2^{j+1}}\right).
\end{equation}
Thus every negative-power antiderivative is finite despite its apparently singular algebraic weight.
\end{corollary}

\begin{proof}
Insert $\nu=-m$ in \eqref{eq:complex-weight-series} and use
$\falling{-m}{j}=(-1)^j\rising{m}{j}$.
\end{proof}

\begin{corollary}[Weights against $\upf$]\label{cor:up-complex-weight}
For $0<x\le1$ and initially $\Re\nu>-1$,
\begin{equation}\label{eq:up-complex-weight}
 \int_0^x t^\nu\upf(t)\dd t
 =\frac{x^{\nu+1}}{\nu+1}-A_\nu(x).
\end{equation}
The right-hand side is the meromorphic continuation in $\nu$, with a single simple pole at $\nu=-1$ of residue $1$.  The regularized integral
\begin{equation}
 \int_0^x t^\nu\bigl(\upf(t)-1\bigr)\dd t=-A_\nu(x)
\end{equation}
continues to an entire function of $\nu$.
\end{corollary}

\newresult Equations \eqref{eq:complex-weight-series} and \eqref{eq:negative-power-series}, including normal convergence for arbitrary complex exponent, were not located explicitly in the audited corpus.  They supply the requested series form for fractional and negative powers without expanding the nowhere-analytic function $F$ in an ordinary Taylor series.

\subsection{Repeated weighted antiderivatives}

\begin{corollary}[Integer-order Volterra primitives with a complex monomial]\label{cor:weighted-nth-primitive}
For $n\in\N$, $0<x\le1$, and $\nu\in\C$,
\begin{equation}\label{eq:weighted-nth-primitive}
 I_0^n[t^\nu F(t)](x)
 =\frac1{(n-1)!}\sum_{r=0}^{n-1}(-1)^r
 \binom{n-1}{r}x^{n-1-r}A_{\nu+r}(x).
\end{equation}
For $\nu=p\in\N_0$, substituting \eqref{eq:integer-weight-antiderivative} gives a finite double sum of dyadically scaled values of $F$; for general $\nu$, substituting \eqref{eq:complex-weight-series} gives a normally convergent series.
\end{corollary}

\begin{proof}
Expand $(x-t)^{n-1}$ in the defining Volterra integral and collect $A_{\nu+r}(x)$.
\end{proof}

\section{Fractional Volterra integration and a two-parameter dyadic recurrence}

\begin{definition}
For $\Re\alpha>0$ and $0<x\le1$, the left Riemann--Liouville integral is
\begin{equation}
 (I_0^\alpha f)(x)=\frac1{\Gamma(\alpha)}
 \int_0^x(x-t)^{\alpha-1}f(t)\dd t.
\end{equation}
\end{definition}

\begin{theorem}[Fractional dilation recurrence]\label{thm:fractional-unweighted}
For $\Re\alpha>1$ and $0<x\le1$,
\begin{equation}\label{eq:fractional-unweighted}
 I_0^\alpha F(x)=2^{\alpha-1}I_0^{\alpha-1}F(x/2).
\end{equation}
More generally, if $m\in\N$ and $m<\Re\alpha$,
\begin{equation}\label{eq:fractional-iterated}
 I_0^\alpha F(x)=
 2^{m\alpha-m(m+1)/2}I_0^{\alpha-m}F(x/2^m).
\end{equation}
At $\alpha=n\in\N$, taking $m=n-1$ recovers \eqref{eq:nth-primitive-F}.
\end{theorem}

\begin{proof}
Substitute $t=2u$ and use $F(2u)=F'(u)/2$:
\[
 I_0^\alpha F(x)=\frac{2^{\alpha-1}}{\Gamma(\alpha)}
 \int_0^{x/2}(x/2-u)^{\alpha-1}F'(u)\dd u.
\]
The boundary terms vanish, and one integration by parts yields \eqref{eq:fractional-unweighted}.  Iteration gives \eqref{eq:fractional-iterated}.
\end{proof}

For a weighted kernel define
\begin{equation}\label{eq:J-alpha-nu}
 J_{\alpha,\nu}(x)=
 \frac1{\Gamma(\alpha)}\int_0^x
 (x-t)^{\alpha-1}t^\nu F(t)\dd t.
\end{equation}
Again, endpoint flatness makes $J_{\alpha,\nu}$ entire in $\nu$ for each fixed $\alpha$ with $\Re\alpha>0$.

\begin{theorem}[Two-parameter fractional recurrence]\label{thm:fractional-weighted}
For $0<x\le1$, $\Re\alpha>1$, and $\nu\in\C$,
\begin{equation}\label{eq:fractional-weighted}
\boxed{
 J_{\alpha,\nu}(x)=2^{\alpha+\nu-1}
 \left[J_{\alpha-1,\nu}(x/2)
       -\nu J_{\alpha,\nu-1}(x/2)\right].}
\end{equation}
The relation reduces the pair $(\alpha,\nu)$ along a two-dimensional dyadic lattice.  When $\nu$ is a nonnegative integer, all branches that lower $\nu$ terminate; when $\alpha$ is an integer, the remaining quantities reduce to \cref{cor:weighted-nth-primitive}.
\end{theorem}

\begin{proof}
After $t=2u$ and $F(2u)=F'(u)/2$,
\[
 J_{\alpha,\nu}(x)=\frac{2^{\alpha+\nu-1}}{\Gamma(\alpha)}
 \int_0^{x/2}(x/2-u)^{\alpha-1}u^\nu F'(u)\dd u.
\]
Integrating by parts differentiates the product of the power and Volterra kernel.  Its two terms are precisely $J_{\alpha-1,\nu}$ and $-\nu J_{\alpha,\nu-1}$.
\end{proof}

\begin{remark}[Fractional derivatives]
For $m-1<\Re\alpha<m$, the Riemann--Liouville derivative
$D_0^\alpha F=D^mI_0^{m-\alpha}F$ can be reduced by \eqref{eq:fractional-unweighted} before differentiation.  Because every derivative of $F$ at $0$ vanishes, Caputo and Riemann--Liouville derivatives coincide on $F$ for positive nonintegral orders.  A useful formalization target is to prove this statement in a function space where the dyadic dilation operator and the fractional Volterra semigroup are both bounded.
\end{remark}

\section{Mellin transforms, logarithmic moments, and Newton interpolation}

\subsection{The entire complex moment function}

\begin{proposition}[Entire moment interpolation]\label{prop:D-entire}
The function
\[
 D(z)=\int_0^1x^z\dd F(x)=\int_0^1x^zF'(x)\dd x
\]
is entire on $\C$.  Its derivatives are
\begin{equation}\label{eq:D-derivatives}
 D^{(m)}(z)=\int_0^1x^z(\log x)^m\dd F(x).
\end{equation}
The integer moments satisfy $d_0=1$ and
\begin{equation}\label{eq:moment-recurrence}
 (2^n-1)d_n=
 \sum_{k=0}^{n-1}\binom nk\frac{d_k}{n-k+1},
 \qquad n\ge1.
\end{equation}
In particular all $d_n$ are rational.
\end{proposition}

\begin{proof}
The density $F'$ is flat at $0$.  On each compact set of $z$-values, a sufficiently high power bound for $F'$ dominates $x^{\Re z}|\log x|^m$ near zero, so differentiation under the integral is valid to every order.  For \eqref{eq:moment-recurrence}, use the distributional fixed point
\[
 X\overset d=\frac{U+X'}2,
\]
where $U$ is uniform on $[0,1]$ and $X'$ is an independent copy of $X$.  Expanding $(U+X')^n$ and isolating the $k=n$ term yields the recurrence.
\end{proof}

The first few moments are
\begin{equation}\label{eq:first-moments}
 1,\quad \frac12,\quad \frac5{18},\quad \frac16,
 \quad \frac{143}{1350},\quad \frac{19}{270},
 \quad \frac{1153}{23814},\ldots
\end{equation}

\subsection{Mellin transforms of $F$ and $\upf$}

\begin{theorem}[Complementary Mellin transforms]\label{thm:mellin-pair}
For $z\in\C$,
\begin{equation}\label{eq:mellin-F}
 \mathcal M_F(z):=\int_0^1x^{z-1}F(x)\dd x
 =\frac{1-D(z)}{z},
\end{equation}
where the quotient has a removable singularity at $z=0$ and is entire.  Initially for $\Re z>0$,
\begin{equation}\label{eq:mellin-up}
 \mathcal M_{\upf,+}(z):=\int_0^1x^{z-1}\upf(x)\dd x
 =\frac{D(z)}{z}.
\end{equation}
The latter identity is its meromorphic continuation to $\C$; it has only one pole, a simple pole at $z=0$ with residue $1$.
\end{theorem}

\begin{proof}
Stieltjes integration by parts gives
\[
 D(z)=1-z\int_0^1x^{z-1}F(x)\dd x
\]
where first justified in a right half-plane and then continued by \cref{prop:D-entire}.  Equation \eqref{eq:mellin-up} follows from $\upf=1-F$ on $[0,1]$.
\end{proof}

\begin{corollary}[Logarithmic integral hierarchy]\label{cor:logarithmic-hierarchy}
For $m\in\N_0$,
\begin{equation}\label{eq:logarithmic-F}
 \int_0^1\frac{(\log x)^m}{x}F(x)\dd x
 =-\frac{D^{(m+1)}(0)}{m+1}.
\end{equation}
Equivalently,
\begin{equation}\label{eq:logarithmic-up-regularized}
 \int_0^1\frac{(\log x)^m}{x}\bigl(\upf(x)-1\bigr)\dd x
 =\frac{D^{(m+1)}(0)}{m+1}.
\end{equation}
In particular,
\begin{equation}\label{eq:ElogX}
 \int_0^1\frac{F(x)}x\dd x=-\E\log X.
\end{equation}
\end{corollary}

\begin{proof}
Differentiate \eqref{eq:mellin-F} at $z=0$ and use the Taylor expansion of $D$.
\end{proof}

\subsection{A Newton series reconstructed from rational moments}

At $x=1$, the dyadic points in \eqref{eq:complex-weight-series} can be eliminated by the repository identity
\begin{equation}\label{eq:dyadic-moment-value}
 d_n=n!\,2^{\binom n2}F(2^{-n}).
\end{equation}
This produces an interpolation formula that reconstructs every complex moment from the rational moments at nonnegative integers.

\begin{corollary}[Newton--Mellin interpolation]\label{cor:newton-mellin}
For every $\nu\in\C$,
\begin{equation}\label{eq:newton-MF}
 \frac{1-D(\nu+1)}{\nu+1}
 =\sum_{j=0}^{\infty}(-1)^j\falling{\nu}{j}
 \frac{d_{j+1}}{(j+1)!},
\end{equation}
with the left side interpreted removably at $\nu=-1$.  Equivalently, for every $z\in\C$,
\begin{equation}\label{eq:newton-D}
\boxed{
 D(z)=1-z\sum_{j=0}^{\infty}(-1)^j
 \falling{z-1}{j}\frac{d_{j+1}}{(j+1)!}.}
\end{equation}
The series converges locally uniformly on $\C$.
\end{corollary}

\begin{proof}
Set $x=1$ in \eqref{eq:complex-weight-series}, use \eqref{eq:dyadic-moment-value}, and combine with \eqref{eq:mellin-F}.  Normal convergence was proved in \cref{thm:complex-weight-series}.  It can also be seen directly: endpoint flatness at $1$ makes $d_n=O_M(n^{-M})$ for every $M$, whereas $\falling{z-1}{j}/j!$ has at most polynomial growth on compact $z$-sets.
\end{proof}

\begin{corollary}[A rapidly convergent logarithmic constant]\label{cor:log-constant-series}
\begin{equation}\label{eq:log-constant-series}
 -\E\log X
 =\int_0^1\frac{F(x)}x\dd x
 =\sum_{n=1}^{\infty}\frac{d_n}{n}
 =0.758239999\ldots .
\end{equation}
The decimal is numerical, while the series identity is exact.
\end{corollary}

\begin{proof}
Put $z=0$ in \eqref{eq:newton-MF}; then
$(-1)^j\falling{-1}{j}=j!$.  The probabilistic identity
$\sum_{n\ge1}d_n/n=\E[-\log(1-X)]=\E[-\log X]$
provides an independent check using symmetry.
\end{proof}

\newresult Formula \eqref{eq:newton-D} is an entire Newton interpolation of a non-elementary complex moment function from rational data.  It creates a direct bridge between the dyadic Fabius values, the moment arithmetic, and complex Mellin analysis.

\subsection{A complex fixed-point equation}

\begin{proposition}[Moment fixed point]\label{prop:D-fixed-point}
For all $z\in\C$, by removable continuation at $z=-1$,
\begin{equation}\label{eq:D-fixed-point}
 D(z)=\frac{2^{-z}}{z+1}
 \E\left[(1+X)^{z+1}-X^{z+1}\right].
\end{equation}
At $z=-1$ this reads
\begin{equation}
 D(-1)=2\E\log\frac{1+X}{X}.
\end{equation}
\end{proposition}

\begin{proof}
Condition on $X'$ in $X\overset d=(U+X')/2$ and integrate $(u+X')^z$ over $0\le u\le1$.
\end{proof}

One may expand $(1+X)^{z+1}$ in a generalized binomial series.  The rapid decay of $d_n$ then gives, initially in a standard convergence region and subsequently by continuation,
\begin{equation}\label{eq:D-binomial-functional}
 D(z)=\frac{2^{-z}}{z+1}
 \left[\sum_{n=0}^{\infty}\binom{z+1}{n}d_n-D(z+1)\right].
\end{equation}
This identity is a promising starting point for analytic difference equations in the Mellin variable.

\subsection{Mellin--Laplace bridges}

Let $B(s)=\E e^{-sX}$, developed in the next section.  For $\Re z<0$,
\begin{equation}\label{eq:mellin-laplace-negative}
 D(z)=\frac1{\Gamma(-z)}\int_0^\infty s^{-z-1}B(s)\dd s.
\end{equation}
More generally, if $m<\Re z<m+1$ with $m\in\N_0$, then
\begin{equation}\label{eq:mellin-laplace-subtracted}
 D(z)=\frac1{\Gamma(-z)}\int_0^\infty s^{-z-1}
 \left[B(s)-\sum_{k=0}^{m}\frac{(-s)^kd_k}{k!}\right]\dd s.
\end{equation}
The subtractions regularize the origin, while $m<\Re z$ controls the polynomial tail at infinity.  These formulas convert asymptotics of the infinite product $B$ into complex-moment information and are suitable for saddle-point analysis of negative moments.

\section{Laplace, Fourier, and Bell--Bernoulli transforms}

\subsection{The one-sided Laplace product}

\begin{theorem}[Laplace product and three exact transforms]\label{thm:laplace-product}
For $s\in\C$,
\begin{equation}\label{eq:B-product}
 B(s)=\E e^{-sX}
 =\prod_{j=1}^{\infty}
 \frac{1-e^{-s/2^j}}{s/2^j},
\end{equation}
with each zero-over-zero factor understood by continuity.  The product is entire and satisfies
\begin{equation}\label{eq:B-dilation}
 B(s)=\frac{1-e^{-s/2}}{s/2}B(s/2).
\end{equation}
Moreover,
\begin{align}
 \int_0^1e^{-sx}\upf(x)\dd x
 &=\frac{1-B(s)}s,\label{eq:laplace-up-one-sided}\\
 \int_0^1e^{-sx}F(x)\dd x
 &=\frac{B(s)-e^{-s}}s,\label{eq:laplace-F-one-sided}\\
 \int_0^1e^{-sG(y)}\dd y
 &=B(s).\label{eq:laplace-inverse}
\end{align}
All quotients are interpreted by removable continuation at $s=0$.
\end{theorem}

\begin{proof}
Independence in \eqref{eq:X-random-series} gives \eqref{eq:B-product}.  The first two integral formulas follow from the master identities in \cref{thm:master-H} with an exponential antiderivative.  The last is the quantile substitution \eqref{eq:quantile-pushforward}.
\end{proof}

For every $p\in\N_0$, differentiation supplies weighted transforms without any new quadrature:
\begin{align}
 \int_0^1x^pe^{-sx}\upf(x)\dd x
 &=(-1)^p\frac{\dd^p}{\dd s^p}\frac{1-B(s)}s,\label{eq:weighted-laplace-up}\\
 \int_0^1x^pe^{-sx}F(x)\dd x
 &=(-1)^p\frac{\dd^p}{\dd s^p}\frac{B(s)-e^{-s}}s,\label{eq:weighted-laplace-F}\\
 \int_0^1G(y)^pe^{-sG(y)}\dd y
 &=(-1)^pB^{(p)}(s).\label{eq:weighted-laplace-G}
\end{align}
For a complex power, Mellin inversion or the series of \cref{thm:complex-weight-series} provides the corresponding continuation.

\subsection{Bernoulli cumulants and Bell-polynomial derivatives}

Let $B_{2m}$ denote the Bernoulli numbers with $B_2=1/6$.  Near the origin,
\begin{equation}\label{eq:log-B-Bernoulli}
 \log B(s)=-\frac{s}{2}
 +\sum_{m=1}^{\infty}
 \frac{B_{2m}}{2m(2m)!(2^{2m}-1)}s^{2m},
 \qquad |s|<4\pi.
\end{equation}
Thus the cumulants of $X$ are
\begin{equation}\label{eq:cumulants-X}
 \kappa_1=\frac12,
 \qquad
 \kappa_{2m}=\frac{B_{2m}}{2m(2^{2m}-1)},
 \qquad
 \kappa_{2m+1}=0\quad(m\ge1).
\end{equation}
If $\mathcal B_n$ is the complete exponential Bell polynomial, then
\begin{equation}\label{eq:moments-Bell}
 d_n=\mathcal B_n(\kappa_1,\ldots,\kappa_n).
\end{equation}
At a general transform parameter,
\begin{equation}\label{eq:B-Bell-derivative}
 B^{(n)}(s)=B(s)\,
 \mathcal B_n\bigl((\log B)'(s),\ldots,(\log B)^{(n)}(s)\bigr).
\end{equation}
Equations \eqref{eq:weighted-laplace-G} and \eqref{eq:B-Bell-derivative} therefore express every exponentially damped inverse moment through Bernoulli data and Bell polynomials.

\subsection{Bilateral transform of $\upf$ and the sinc product}

Since $Y=2X-1$ has density $\upf$,
\begin{equation}\label{eq:bilateral-up-laplace}
 \int_{-1}^{1}e^{-su}\upf(u)\dd u
 =e^sB(2s)
 =\prod_{j=1}^{\infty}
 \frac{\sinh(s/2^j)}{s/2^j}.
\end{equation}
With the angular-frequency convention
\[
 \phi_Y(\omega)=\int_{-1}^{1}e^{\ii\omega u}\upf(u)\dd u,
\]
one obtains
\begin{equation}\label{eq:angular-sinc-product}
 \phi_Y(\omega)=\prod_{j=1}^{\infty}
 \sinc(\omega/2^j),
 \qquad \sinc z=\frac{\sin z}{z}.
\end{equation}
In the repository convention
$\widehat f(\xi)=\int f(x)e^{-2\pi\ii\xi x}\dd x$,
this becomes
\begin{equation}
 \widehat{\upf}(\xi)=
 \prod_{n=0}^{\infty}\sinc(\pi\xi/2^n).
\end{equation}
Polynomially weighted Fourier transforms follow from
\begin{equation}\label{eq:weighted-fourier}
 \int_{-1}^{1}u^ne^{\ii\omega u}\upf(u)\dd u
 =\ii^{-n}\phi_Y^{(n)}(\omega).
\end{equation}
Away from the zeros of the product, $\phi_Y^{(n)}$ can again be organized by Bell polynomials applied to the logarithmic derivatives of the sinc factors.  At a zero, the entire derivative in \eqref{eq:weighted-fourier}, rather than a singular logarithmic representation, is the correct object.

\section{Integral calculus of the inverse Fabius function}

\subsection{Global inverse moments and local primitives}

The quantile push-forward immediately turns every inverse-Fabius moment into a moment of $X$.

\begin{theorem}[Entire inverse-power transform]\label{thm:inverse-entire-moments}
For every $z\in\C$,
\begin{equation}\label{eq:inverse-entire-moments}
 \int_0^1G(y)^z\dd y=D(z).
\end{equation}
The integral converges for all complex $z$, including arbitrarily large negative real part.
\end{theorem}

\begin{proof}
Equation \eqref{eq:inverse-entire-moments} is \eqref{eq:quantile-pushforward} with $\Phi(x)=x^z$.  Convergence follows from \cref{prop:D-entire}; equivalently, $G(y)$ tends to zero more slowly than every positive power of $y$, so negative powers of $G$ remain integrable at $y=0$.
\end{proof}

\begin{theorem}[Exact local inverse-power primitive]\label{thm:inverse-local-primitive}
Let $0\le y\le1$, put $x=G(y)$, and let $p\in\C$.  Then
\begin{equation}\label{eq:inverse-local-primitive}
\boxed{
 \int_0^yG(t)^p\dd t
 =y x^p-pA_{p-1}(x).}
\end{equation}
For $p\in\N$, \cref{thm:integer-weight-antiderivative} makes the right side a finite sum over $F(x/2^j)$, $1\le j\le p$.  For arbitrary complex $p$, \cref{thm:complex-weight-series} gives a normally convergent dyadic series.
\end{theorem}

\begin{proof}
Use $t=F(u)$, so $u=G(t)$ and $\dd t=F'(u)\dd u$:
\[
 \int_0^yG(t)^p\dd t
 =\int_0^xu^pF'(u)\dd u
 =yx^p-p\int_0^xu^{p-1}F(u)\dd u.
\]
The lower boundary vanishes for every complex $p$ by endpoint flatness.
\end{proof}

For example,
\begin{align}
 \int_0^yG(t)\dd t
 &=yG(y)-F(G(y)/2),\label{eq:inverse-primitive-p1}\\
 \int_0^yG(t)^2\dd t
 &=yG(y)^2-2G(y)F(G(y)/2)+4F(G(y)/4),\label{eq:inverse-primitive-p2}\\
 \int_0^yG(t)^3\dd t
 &=yG(y)^3-3G(y)^2F(G(y)/2)
 +12G(y)F(G(y)/4)-48F(G(y)/8).
\end{align}
At $y=1$, these identities reduce to $\int_0^1G^p=d_p$.

\begin{corollary}[A weighted CDF--quantile duality]\label{cor:weighted-quantile-duality}
For $p\in\C$ and $\Re q>0$,
\begin{equation}\label{eq:weighted-quantile-duality}
 \int_0^1G(y)^p y^{q-1}\dd y
 =\frac1q-\frac pq\int_0^1x^{p-1}F(x)^q\dd x.
\end{equation}
The endpoint at $x=0$ causes no restriction on $p$ because $F(x)^q$ is flat whenever $\Re q>0$.
\end{corollary}

\begin{proof}
Put $y=F(x)$ and integrate $x^p\dd(F(x)^q)/q$ by parts.
\end{proof}

A second universal identity, useful for inverse derivatives, is
\begin{equation}\label{eq:Gprime-pushforward}
 \int_0^1\Phi(G(y))G'(y)\dd y=\int_0^1\Phi(x)\dd x
\end{equation}
whenever either side is absolutely integrable.

\subsection{All derivatives of the inverse}

Set $f=F'$.  Define differential polynomials $Q_n[f]$ recursively by
\begin{equation}\label{eq:Q-recurrence}
 Q_1=1,
 \qquad
 Q_{n+1}=fQ_n'-(2n-1)f'Q_n.
\end{equation}

\begin{theorem}[All-order inverse derivative formula]\label{thm:inverse-all-derivatives}
If $x=G(y)$, then for every $n\ge1$,
\begin{equation}\label{eq:inverse-all-derivatives}
\boxed{
 G^{(n)}(y)=\frac{Q_n[f](x)}{f(x)^{2n-1}}.}
\end{equation}
The first cases are
\begin{align}
 G'(y)&=\frac1{F'(x)},\notag\\
 G''(y)&=-\frac{F''(x)}{F'(x)^3},\notag\\
 G'''(y)&=\frac{3F''(x)^2-F'(x)F'''(x)}{F'(x)^5},\label{eq:inverse-first-derivatives}\\
 G^{(4)}(y)&=\frac{-15F''(x)^3+10F'(x)F''(x)F'''(x)
 -(F'(x))^2F^{(4)}(x)}{F'(x)^7}.
\end{align}
Using the signed extension,
\begin{equation}\label{eq:forward-derivative-dilation}
 F^{(m)}(x)=2^{\binom{m+1}{2}}\mathcal F(2^mx),
 \qquad 0<x<1,
\end{equation}
so every inverse derivative is a rational expression in finitely many dyadically dilated values of $\mathcal F$.
\end{theorem}

\begin{proof}
The differentiation operator in the inverse coordinate is
$\frac{\dd}{\dd y}=f(x)^{-1}\frac{\dd}{\dd x}$.  Assuming \eqref{eq:inverse-all-derivatives} at order $n$ and differentiating gives
\[
 \frac1f\frac{\dd}{\dd x}\left(Q_nf^{-(2n-1)}\right)
 =\frac{fQ_n'-(2n-1)f'Q_n}{f^{2n+1}},
\]
which is exactly \eqref{eq:Q-recurrence}.
\end{proof}

The recurrence is a compact alternative to a full Fa\`a di Bruno expansion.  The $Q_n$ are closely related to the inverse-function Bell polynomials; \eqref{eq:Q-recurrence} is often more efficient computationally because it preserves the natural denominator $f^{2n-1}$.

\subsection{Sharp integrability of the first inverse derivative}

\begin{theorem}[$L^q$ threshold for $G'$]\label{thm:Gprime-Lq}
For $q>0$,
\begin{equation}\label{eq:Gprime-Lq-change}
 \int_0^1|G'(y)|^q\dd y
 =\int_0^1F'(x)^{1-q}\dd x.
\end{equation}
Consequently,
\begin{equation}\label{eq:Gprime-Lq-threshold}
 G'\in L^q(0,1)
 \quad\Longleftrightarrow\quad 0<q\le1.
\end{equation}
In particular $G\in W^{1,1}(0,1)$ but $G\notin W^{1,q}(0,1)$ for any $q>1$.
\end{theorem}

\begin{proof}
Substitute $y=F(x)$ in the left side.  If $0<q\le1$, the exponent $1-q$ is nonnegative and $F'$ is bounded.  If $q>1$, flatness gives $F'(x)\le C_Nx^N$ near $0$ for every $N$.  Choose $N$ with $N(q-1)\ge1$; then $F'(x)^{1-q}$ dominates a nonintegrable multiple of $x^{-N(q-1)}$.
\end{proof}

The complementary density--quantile identity is
\begin{equation}\label{eq:inverse-density-power}
 \int_0^1G'(y)^{-q}\dd y
 =\int_0^1F'(x)^{q+1}\dd x,
\end{equation}
whenever finite.  Two notable specializations are
\begin{align}
 \int_0^1\log G'(y)\dd y
 &=-\int_0^1F'(x)\log F'(x)\dd x\notag\\
 &=-\int_{-1}^{1}\upf(u)\log\bigl(2\upf(u)\bigr)\dd u,
 \label{eq:quantile-entropy}\\
 \int_0^1\upf(2G(y)-1)^r\dd y
 &=\int_{-1}^{1}\upf(u)^{r+1}\dd u,
 \qquad r>-1.\label{eq:inverse-up-power}
\end{align}
Formula \eqref{eq:quantile-entropy} identifies the mean logarithmic slope of the quantile with the differential entropy of the Fabius density.

\section{Endpoint asymptotics of inverse primitives}

The audited inverse-function reports derive a Lambert--$W$ endpoint law with a small periodic fluctuation.  Write
\begin{equation}\label{eq:rho-definition}
 T=\log(1/y),\qquad L=\log2,
 \qquad \rho=\sqrt{\frac{2T}{L}},
 \qquad \beta=\frac1L+\frac12.
\end{equation}
Their leading inversion gives
\begin{equation}\label{eq:G-leading-endpoint}
 G(y)\sim \frac{\sqrt T}{e\sqrt L}
 \exp\!\left(-\sqrt{2LT}\right),
 \qquad y\downarrow0.
\end{equation}
In particular $G$ is slowly varying at zero: $G(cy)/G(y)\to1$ for each fixed $c>0$.

Let $\Psi$ denote the bounded period-$1$ phase function in the repository's refined inverse asymptotic and define
\begin{equation}\label{eq:a-rho}
 a(\rho)=\frac{\Psi'(\rho+\beta)-1}{L}.
\end{equation}
The logarithmic elasticity is
\begin{equation}\label{eq:epsilon-expansion}
 \epsilon(y):=\frac{yG'(y)}{G(y)}
 =\frac1\rho+\frac{a(\rho)}{\rho^2}
 +O\!\left(\frac{\log\rho}{\rho^3}\right).
\end{equation}

\begin{theorem}[Refined inverse-power primitive]\label{thm:inverse-primitive-asymptotic}
Assume the derivative-controlled form of the all-orders inverse phase expansion underlying \eqref{eq:epsilon-expansion}.  For each fixed $p\in\C$, as $y\downarrow0$,
\begin{align}
 A_{p-1}(G(y))
 &=yG(y)^p\left[
 \frac1\rho+\frac{a(\rho)-p}{\rho^2}
 +O\!\left(\frac{\log\rho}{\rho^3}\right)
 \right],\label{eq:A-at-G-asymptotic}\\
 \int_0^yG(t)^p\dd t
 &=yG(y)^p\left[
 1-\frac p\rho+\frac{p\bigl(p-a(\rho)\bigr)}{\rho^2}
 +O\!\left(\frac{\log\rho}{\rho^3}\right)
 \right].\label{eq:inverse-primitive-refined}
\end{align}
In particular,
\begin{equation}\label{eq:F-half-G-asymptotic}
 F(G(y)/2)=yG(y)\left[
 \frac1\rho+\frac{a(\rho)-1}{\rho^2}
 +O\!\left(\frac{\log\rho}{\rho^3}\right)
 \right].
\end{equation}
The period-$1$ ripple survives integration, but first appears one order below the leading correction.
\end{theorem}

\begin{proof}
Putting $u=G(t)$ in $A_{p-1}(G(y))$ gives the exact identity
\begin{equation}\label{eq:A-G-epsilon}
 A_{p-1}(G(y))=\int_0^yG(t)^p\epsilon(t)\dd t.
\end{equation}
Let $y=e^{-T}$ and $H(T)=G(e^{-T})^p\epsilon(e^{-T})$.  Then
\[
 A_{p-1}(G(y))=\int_T^\infty e^{-u}H(u)\dd u
 =e^{-T}\bigl(H(T)+H'(T)+O(H''(T))\bigr).
\]
The derivative-controlled version of \eqref{eq:epsilon-expansion} gives
\[
 H'(T)=G(e^{-T})^p\left[-\frac p{\rho^2}
 +O\!\left(\frac{\log\rho}{\rho^3}\right)\right],
 \qquad
 H''(T)=O\!\left(G(e^{-T})^p\frac{\log\rho}{\rho^3}\right).
\]
Substitution proves \eqref{eq:A-at-G-asymptotic}.  Formula \eqref{eq:inverse-primitive-refined} follows from \eqref{eq:inverse-local-primitive}, and \eqref{eq:F-half-G-asymptotic} is the case $p=1$ because $A_0(x)=F(x/2)$.
\end{proof}

Without the second-order phase information, Karamata's theorem already gives the robust leading form
\begin{equation}\label{eq:inverse-primitive-Karamata}
 \int_0^y t^{q-1}G(t)^p\dd t
 \sim\frac{y^qG(y)^p}{q},
 \qquad q>0,
\end{equation}
for every fixed real $p$.  Equation \eqref{eq:inverse-primitive-refined} resolves the first two corrections in the important case $q=1$.

\section{Stieltjes, Cauchy, and logarithmic transforms}\label{sec:stieltjes}

For $w\in\C\setminus[0,1]$ define the Cauchy--Stieltjes transform
\begin{equation}\label{eq:S-def}
 S(w)=\E\frac1{w-X}
 =\int_0^1\frac{\dd F(x)}{w-x}.
\end{equation}

\begin{theorem}[Four equivalent Stieltjes representations]\label{thm:stieltjes-representations}
For $w\notin[0,1]$,
\begin{align}
 S(w)
 &=\int_0^1\frac{\dd y}{w-G(y)},\label{eq:S-quantile}\\
 &=\frac1w+\int_0^1\frac{\upf(t)}{(w-t)^2}\dd t,\label{eq:S-up}\\
 &=\frac1{w-1}-\int_0^1\frac{F(t)}{(w-t)^2}\dd t.\label{eq:S-F}
\end{align}
For $|w|>1$,
\begin{equation}\label{eq:S-moment-series}
 S(w)=\sum_{n=0}^{\infty}\frac{d_n}{w^{n+1}}.
\end{equation}
The transform of Rvachev's centered density is
\begin{equation}\label{eq:Cauchy-up}
 \int_{-1}^{1}\frac{\upf(u)}{w-u}\dd u
 =\frac12S\!\left(\frac{w+1}{2}\right),
 \qquad w\notin[-1,1].
\end{equation}
\end{theorem}

\begin{proof}
Equation \eqref{eq:S-quantile} is the quantile push-forward.  Apply \cref{thm:master-H} to $H(t)=(w-t)^{-1}$ to obtain \eqref{eq:S-up}; Stieltjes integration by parts gives \eqref{eq:S-F}.  Geometric expansion gives \eqref{eq:S-moment-series}.  Finally use $u=2x-1$ and $\dd F(x)=2\upf(2x-1)\dd x$.
\end{proof}

A primitive in the spectral variable produces logarithmic kernels.

\begin{corollary}[Logarithmic potentials]\label{cor:logarithmic-potentials}
Choose a branch of the logarithm on a simply connected domain avoiding $[0,1]$.  Then
\begin{align}
 \int_0^1\frac{F(x)}{w-x}\dd x
 &=\E\log(w-X)-\log(w-1),\label{eq:log-potential-F}\\
 \int_0^1\frac{\upf(x)}{w-x}\dd x
 &=\log w-\E\log(w-X),\label{eq:log-potential-up}\\
 \int_0^1\log(w-G(y))\dd y
 &=\E\log(w-X).\label{eq:log-potential-G}
\end{align}
Differentiating any one of these equations recovers the corresponding Stieltjes transform.
\end{corollary}

\begin{proof}
Integrate by parts using $\frac{\dd}{\dd x}\log(w-x)=-(w-x)^{-1}$ and apply the quantile push-forward to the last formula.
\end{proof}

Boundary values of $S$ encode the Fabius density through the Sokhotski--Plemelj formula.  Thus the Cauchy transform packages moment arithmetic outside the interval and reconstructs $F'$ on the cut.  Its continued-fraction expansion, determined by the Hankel determinants of $(d_n)$, is developed below.

\section{Products, order statistics, and autocorrelation}\label{sec:products-correlations}

\subsection{Products of $F$ and $1-F$}

For positive integers $m,n$, let $A_m=\max(X_1,\ldots,X_m)$ and $B_n=\min(X'_1,\ldots,X'_n)$, with the two samples independent.  Then
\begin{equation}\label{eq:order-statistic-product}
 \int_0^1F(x)^m\bigl(1-F(x)\bigr)^n\dd x
 =\E(B_n-A_m)_+.
\end{equation}
Indeed, the integrand is $\Prob(A_m\le x<B_n)$, and integration measures the random interval length.  Two useful special cases are
\begin{equation}\label{eq:powers-order-statistics}
 \int_0^1F(x)^m\dd x=1-\E A_m,
 \qquad
 \int_0^1\upf(x)^m\dd x=\E B_m.
\end{equation}
Symmetry gives equality of the two integrals.

The first mixed product admits several especially transparent forms.

\begin{theorem}[A Fabius Gini integral]\label{thm:Fabius-Gini}
Let $X'$ be an independent copy of $X$.  Then
\begin{align}
 C_F&:=\int_0^1F(x)\bigl(1-F(x)\bigr)\dd x
 =\frac12\E|X-X'|,\label{eq:Gini-probability}\\
 &=\frac1\pi\int_0^\infty
 \frac{1-|\phi_X(t)|^2}{t^2}\dd t,\label{eq:Gini-Fourier}
\end{align}
where
\begin{equation}\label{eq:phi-X}
 \phi_X(t)=e^{\ii t/2}
 \prod_{j=1}^{\infty}\sinc(t/2^{j+1}).
\end{equation}
Equivalently,
\begin{equation}\label{eq:Gini-centered-Fourier}
 C_F=\frac1{2\pi}\int_0^\infty
 \frac{1-\phi_Y(t)^2}{t^2}\dd t.
\end{equation}
Furthermore,
\begin{equation}\label{eq:Gini-up-L2}
 \int_{-1}^{1}\upf(u)^2\dd u=1-2C_F,
 \qquad
 \int_0^1F(x)^2\dd x=\frac12-C_F.
\end{equation}
Numerically,
\begin{equation}\label{eq:Gini-numerical}
 C_F\approx0.0955632.
\end{equation}
\end{theorem}

\begin{proof}
The standard one-dimensional Gini identity gives
$\E|X-X'|=2\int F(1-F)$.  The Fourier representation follows from
$|z|=\frac2\pi\int_0^\infty(1-\cos tz)t^{-2}\dd t$
and the characteristic function of $X-X'$.  Since $X=(Y+1)/2$, \eqref{eq:Gini-centered-Fourier} follows by scaling.  Finally, on $[0,1]$ one has $\upf=1-F$, and symmetry gives
$\int F^2=\int(1-F)^2$.  Expanding $F(1-F)$ proves \eqref{eq:Gini-up-L2}.
\end{proof}

\subsection{The autocorrelation transform}

Define
\begin{equation}\label{eq:up-autocorrelation}
 R(a)=\int_{\R}\upf(x)\upf(x-a)\dd x.
\end{equation}
Because $\upf$ is even, $R=\upf*\upf$; it is an even $C^\infty$ probability density supported on $[-2,2]$.  In angular-frequency convention,
\begin{equation}\label{eq:up-autocorrelation-Fourier}
 R(a)=\frac1{2\pi}\int_{\R}e^{-\ii\omega a}
 \prod_{j=1}^{\infty}\sinc^2(\omega/2^j)\dd\omega.
\end{equation}
If $\mu_n=\int u^n\upf(u)\dd u$, then
\begin{equation}\label{eq:autocorrelation-moments}
 \int_{-2}^{2}a^nR(a)\dd a
 =\sum_{k=0}^n\binom nk\mu_k\mu_{n-k}.
\end{equation}
Only even indices survive.  More generally, correlations of derivatives are obtained by multiplying the integrand in \eqref{eq:up-autocorrelation-Fourier} by powers of $\ii\omega$.  These identities translate questions about products of shifted Rvachev functions into the arithmetic of a squared lacunary sinc product.

\section{Numerical experiments}

\subsection{Finite Thue--Morse spline model}

The accompanying program uses the finite random sum
\begin{equation}
 Y_n=\sum_{j=1}^n\frac{V_j}{2^j}
\end{equation}
and its compactly supported spline density.  Let
$\tau_j=(-1)^{s_2(j)}$, $A_n=1-2^{-n}$, and $(r)_+=\max(r,0)$.  Then
\begin{align}
 p_n(y)&=\frac{2^{\binom n2}}{(n-1)!}
 \sum_{j=0}^{2^n-1}\tau_j
 \left(y+A_n-\frac{j}{2^{n-1}}\right)_+^{n-1},\label{eq:finite-density}\\
 P_n(y)&=\frac{2^{\binom n2}}{n!}
 \sum_{j=0}^{2^n-1}\tau_j
 \left(y+A_n-\frac{j}{2^{n-1}}\right)_+^{n}.
 \label{eq:finite-cdf}
\end{align}
The finite Fabius CDF is $F_n(x)=P_n(2x-1)$.  Long-double arithmetic and symmetry are used to suppress cancellation.  Exact rational moments are generated independently from \eqref{eq:moment-recurrence}.

\subsection{Results and reproducibility}

The first table gives the largest residual among several test points for
$\int_0^xF_n(t)\dd t\approx F_n(x/2)$, together with the finite-rank approximation to $C_F$.  The second tests \eqref{eq:integer-weight-antiderivative} at $x=0.83$ and rank $10$.  The third shows convergence of \eqref{eq:newton-D} at $z=0.5+0.7\ii$.

\begin{center}
\begin{tabular}{rcc}
\toprule
Spline rank & max. primitive residual & $\int_0^1F_n(1-F_n)$\\
\midrule
6 & $1.873e-05$ & $0.095552260300$\\
8 & $1.170e-06$ & $0.095562546323$\\
10 & $7.313e-08$ & $0.095563189160$\\
\bottomrule
\end{tabular}
\par\medskip
\begin{tabular}{rccc}
\toprule
$p$ & quadrature & dyadic formula & difference\\
\midrule
0 & $0.330437473102$ & $0.330437457904$ & $1.520e-08$\\
1 & $0.205934882125$ & $0.205935065265$ & $-1.831e-07$\\
2 & $0.135359629291$ & $0.135359392297$ & $2.370e-07$\\
3 & $0.092437619662$ & $0.092438027029$ & $-4.074e-07$\\
5 & $0.046620893512$ & $0.046622072851$ & $-1.179e-06$\\
\bottomrule
\end{tabular}
\par\medskip
\begin{tabular}{rcc}
\toprule
Newton terms & approximation error & approximation\\
\midrule
5 & $8.605e-03$ & $0.601640735-0.321223181\,i$\\
10 & $1.041e-03$ & $0.598727644-0.314110902\,i$\\
20 & $6.790e-05$ & $0.598823615-0.313135185\,i$\\
40 & $2.260e-06$ & $0.598859558-0.313079898\,i$\\
60 & $2.590e-07$ & $0.598861332-0.313078830\,i$\\
100 & $1.245e-07$ & $0.598861546-0.313078775\,i$\\
\bottomrule
\end{tabular}

\end{center}

The beta integral \eqref{eq:beta-family} was also evaluated directly for several nonintegral exponent pairs; residuals were between $4\times10^{-16}$ and $6\times10^{-14}$.  A rank-$10$, 400-node quantile quadrature checked \eqref{eq:laplace-inverse} with errors
\begin{equation}
 5.7\times10^{-9},\quad3.6\times10^{-8},\quad7.3\times10^{-8}
\end{equation}
for $s=0.7,2.5,5$, respectively.  The partial sums
$\sum_{n\le200}d_n/n$ give $0.758239999524$, agreeing with \eqref{eq:log-constant-series}.

The finite spline is not forced to obey the infinite-rank dyadic primitive identity exactly; its residual falls by about a factor of $16$ whenever the rank increases by $2$, consistent with convergence to the limiting law.  By contrast, the beta substitution is exact for every smooth finite-rank CDF, so its residual is essentially floating-point roundoff.

The archive contains \texttt{numerical\_experiments.py}.  Running
\begin{lstlisting}[language=bash]
python numerical_experiments.py
\end{lstlisting}
regenerates \texttt{numerical\_results.txt} and
\texttt{numerical\_results.tex}.  The code documents every normalization, the finite spline formula, quadrature, bisection inversion, and rational-moment recurrence.

\section{An arithmetic continued fraction generated by the moments}

The Stieltjes transform admits one more exact structure that appears not to have been exploited in the audited corpus.  Let $P_n$ be the monic polynomials orthogonal with respect to $\dd F$ on $[0,1]$, normalized by $P_{-1}=0$, $P_0=1$.

\begin{proposition}[Rational Jacobi parameters]\label{prop:rational-Jacobi}
There are positive rational numbers $b_n$ such that
\begin{equation}\label{eq:orthogonal-recurrence}
 P_{n+1}(x)=\left(x-\frac12\right)P_n(x)-b_nP_{n-1}(x),
 \qquad n\ge0,
\end{equation}
where the term containing $b_0$ is absent.  In particular,
\begin{equation}\label{eq:first-Jacobi}
 b_1=\frac1{36},
 \qquad b_2=\frac8{225}.
\end{equation}
Consequently the Cauchy transform has the Stieltjes $J$-fraction
\begin{equation}\label{eq:S-J-fraction}
 S(w)=\cfrac{1}{w-\frac12-
 \cfrac{b_1}{w-\frac12-
 \cfrac{b_2}{w-\frac12-\ddots}}}.
\end{equation}
Every finite convergent is a rational function with rational coefficients.
\end{proposition}

\begin{proof}
Gram--Schmidt applied to $1,x,x^2,\ldots$ uses only the rational moments $d_n$, so every coefficient of $P_n$ and every squared norm
$h_n=\int P_n^2\dd F$ is rational.  Positivity of the measure gives $h_n>0$ and $b_n=h_n/h_{n-1}>0$.  Symmetry under $x\mapsto1-x$ implies
$P_n(1-x)=(-1)^nP_n(x)$, hence
$\int xP_n(x)^2\dd F(x)=h_n/2$ and every diagonal Jacobi coefficient equals $1/2$.  The standard resolvent expansion of the Jacobi matrix gives \eqref{eq:S-J-fraction}.  Finally \(b_1=\mathrm{Var}(X)=1/36\); using the fourth centered moment $19/10800$ gives $b_2=8/225$.
\end{proof}

This rational Jacobi matrix is a new arithmetic object attached to the Fabius law.  Determining valuations, denominator growth, and the rate at which $b_n$ approaches the universal interval value is a natural continuation of the repository's dyadic arithmetic program.

\section{Conjectures and frontier directions}

\subsection{Higher inverse derivatives}

Let
\begin{equation}
 R_n(y)=\frac{y^nG^{(n)}(y)}{G(y)},
 \qquad \delta=y\frac{\dd}{\dd y}.
\end{equation}
Direct differentiation gives the exact recurrence
\begin{equation}\label{eq:Rn-recurrence}
 R_1=\epsilon,
 \qquad
 R_{n+1}=(\epsilon-n)R_n+\delta R_n.
\end{equation}
The first-order elasticity and the all-orders Lambert phase strongly suggest the following.

\begin{conjecture}[All-order inverse endpoint law]\label{conj:inverse-derivatives}
For every fixed $n\ge1$,
\begin{equation}\label{eq:inverse-derivative-conjecture}
 G^{(n)}(y)\sim(-1)^{n-1}(n-1)!
 \frac{G(y)}{y^n\rho(y)},
 \qquad y\downarrow0.
\end{equation}
Consequently, for $n\ge1$ and $q>0$,
\begin{equation}\label{eq:higher-inverse-Lq-conjecture}
 G^{(n)}\in L^q(0,1)
 \quad\Longleftrightarrow\quad 0<q\le\frac1n.
\end{equation}
The case $n=1$ is \cref{thm:Gprime-Lq}.  At the borderline $q=1/n$, the residual factor $G(y)^{1/n}/\rho(y)^{1/n}$ should make the logarithmic-coordinate integral convergent.
\end{conjecture}

A proof should follow by establishing derivative-stable error estimates for the all-orders inverse phase, then inducting in \eqref{eq:Rn-recurrence}.  This is more delicate than formal differentiation because an undifferentiated asymptotic remainder does not control its derivatives.

\subsection{Negative moments and a dual endpoint saddle}

The entire function $D$ grows rapidly on the negative real axis.  The dyadic endpoint formula, Stirling's approximation, and the forward saddle suggest the coarse scale
\begin{conjecture}[Negative-moment saddle]\label{conj:negative-moments}
As $r\to+\infty$,
\begin{equation}\label{eq:negative-moment-conjecture}
 \log D(-r)=\frac{\log2}{2}r^2-r\log r+O(r).
\end{equation}
After subtracting the complete smooth expansion in $r$, a bounded period-$1$ function of the corrected phase
$r-\log_2r+O(1)$ should remain.
\end{conjecture}

This problem is a Mellin-dual counterpart of the small-$x$ and inverse-Fabius endpoint analysis.  Equation \eqref{eq:mellin-laplace-negative} offers a second route through the large-$s$ saddle of $B(s)$.

\subsection{Autocorrelation arithmetic}

The squared sinc product in \eqref{eq:up-autocorrelation-Fourier} has a refinement mask obtained by squaring the original mask.  It is reasonable to expect a dyadic arithmetic theory for $R(a)$ analogous to that of $\upf(a)$, but with triangular rather than uniform digit laws.

\begin{conjecture}[Dyadic autocorrelation arithmetic]\label{conj:autocorrelation-rational}
For every dyadic rational $a\in[-2,2]$, the value $R(a)$ belongs to the field generated by the Fabius dyadic values and, in fact, is rational after the natural power-of-two normalization determined by the spline rank.  Its odd-prime denominators should be governed by products of factors $2^{2m}-1$ with multiplicities reflecting the squared mask.
\end{conjecture}

Even a weaker exact recurrence for $R(k/2^m)$ would be useful: it would turn shifted products of $\upf$ into arithmetic data and would provide certified values of $\|\upf\|_2^2=R(0)$.

\subsection{Fractional operator spectrum}

Equation \eqref{eq:fractional-weighted} suggests studying the dilation--Volterra operator algebra generated by
\[
 (Df)(x)=f(x/2),
 \qquad
 (V_\alpha f)(x)=\frac1{\Gamma(\alpha)}
 \int_0^x(x-t)^{\alpha-1}f(t)\dd t.
\]
The Fabius function is a cyclic vector for a noncommutative relation between these operators.  Promising questions include:
\begin{itemize}
\item meromorphic continuation of $J_{\alpha,\nu}$ jointly in $(\alpha,\nu)$ after canonical endpoint subtractions;
\item spectral decompositions of the dilation generator in Mellin coordinates;
\item base-$b$ analogues, replacing the dyadic triangular exponents by $b$-adic quadratic forms;
\item a fractional analogue of the signed derivative hierarchy \eqref{eq:signed-differential}.
\end{itemize}

\subsection{Mellin zeros, growth, and acceleration}

The Newton series \eqref{eq:newton-D} can be evaluated entirely from rational moments.  Its convergence is superalgebraic but not generally geometric.  The endpoint phase suggests splitting $d_n$ into a smooth saddle part and a tiny periodic residual, then applying Euler or N\o rlund acceleration to the Newton basis.  This may reveal:
\begin{itemize}
\item the order and indicator function of the entire function $D(z)$;
\item zero patterns of $D$, $1-D$, and the regularized Mellin transforms;
\item functional relations between those zeros and the zeros of the Laplace product $B$;
\item certified high-precision negative moments, inaccessible to direct quadrature.
\end{itemize}

\subsection{Orthogonal-polynomial arithmetic}

The rational coefficients $b_n$ in \eqref{eq:orthogonal-recurrence} deserve a separate arithmetic study.  General orthogonal-polynomial theory predicts interval-universal limiting behavior because the density is positive on $(0,1)$, but the superflat endpoints may leave a slow, log-periodic signature in the convergence rate.  Concrete problems are:
\begin{itemize}
\item determine $2$-adic and odd-prime valuations of $b_n$;
\item derive a nonlinear recurrence directly from the moment recurrence \eqref{eq:moment-recurrence};
\item identify the first log-periodic term in the large-$n$ Jacobi asymptotics;
\item connect Pad\'e approximants of $S$ to exact dyadic evaluations of $F$.
\end{itemize}

\section{Novelty ledger and formalization roadmap}

\begin{longtable}{@{}p{0.28\textwidth}p{0.18\textwidth}p{0.43\textwidth}@{}}
\toprule
Result & Status in audit & Main proof mechanism\\
\midrule
\endhead
$\int_0^xF=F(x/2)$ and integer iterates
& baseline, partly implicit
& Fabius differential dilation and normalization\\
\addlinespace
Finite formula \eqref{eq:integer-weight-antiderivative}
& not located explicitly
& repeated integration by parts with shrinking argument\\
\addlinespace
Complex series \eqref{eq:complex-weight-series}
& new relative to audit
& exact recurrence plus endpoint-flat factorial bound\\
\addlinespace
Newton--Mellin formula \eqref{eq:newton-D}
& new relative to audit
& dyadic endpoint/moment identity at $x=1$\\
\addlinespace
Fractional recurrence \eqref{eq:fractional-weighted}
& new relative to audit
& Riemann--Liouville kernel and one integration by parts\\
\addlinespace
Local inverse primitive \eqref{eq:inverse-local-primitive}
& not located explicitly
& quantile substitution followed by Stieltjes integration by parts\\
\addlinespace
All inverse derivatives \eqref{eq:inverse-all-derivatives}
& standard inverse calculus, new packaging
& differential-polynomial recurrence preserving $F'^{2n-1}$\\
\addlinespace
Sharp $L^q$ threshold \eqref{eq:Gprime-Lq-threshold}
& not located explicitly
& change of variables and infinite endpoint flatness\\
\addlinespace
Refined inverse primitive \eqref{eq:inverse-primitive-refined}
& new consequence of repository asymptotics
& elasticity identity and Watson expansion\\
\addlinespace
Stieltjes/logarithmic transforms
& new systematic synthesis
& master CDF--survival and quantile identities\\
\addlinespace
Gini/autocorrelation formulas
& new synthesis
& order statistics and Fourier representation of absolute distance\\
\addlinespace
Rational Jacobi parameters
& not located explicitly
& rational moments, symmetry, and Gram--Schmidt\\
\bottomrule
\end{longtable}

The following sequence is realistic for Lean formalization:
\begin{enumerate}
\item formalize $A_p$ for $p\in\N_0$, prove \eqref{eq:A-recurrence}, and derive the finite sum by induction;
\item formalize the $n$th Volterra primitive \eqref{eq:weighted-nth-primitive} using finite binomial expansion;
\item define $D(z)$ first on real exponents, prove the moment recurrence and Mellin identities, then add complex analyticity;
\item prove the finite Newton identities at integer $z$ before introducing normal convergence of \eqref{eq:newton-D};
\item formalize the inverse local primitive and the $Q_n$ recurrence, which require only one-dimensional calculus;
\item treat $L^q$ divergence using the existing endpoint-flat lemmas;
\item postpone the refined periodic asymptotics until derivative-controlled error bounds are available in the formal library.
\end{enumerate}

\section{Conclusion}

The Fabius--Rvachev system is unusually well suited to integration: differentiation doubles the argument, while normalized integration halves it.  Once this elementary observation is combined with endpoint flatness, it yields more than the familiar primitive $F(x/2)$.  Polynomial weights produce finite dyadic sums; arbitrary complex weights produce normally convergent dyadic series; the endpoint values turn those series into an entire Newton interpolation of complex moments; and quantile substitution transfers the same structure to the inverse Fabius function.

The transform picture is similarly unified.  Mellin transforms are controlled by the entire moment function $D$, Laplace transforms by the product $B$, Fourier transforms by the lacunary sinc product, and Stieltjes transforms by a rational-moment Jacobi matrix.  Products of $F$ become order-statistic gaps, while shifted products of $\upf$ become autocorrelations of the squared sinc product.  The resulting formulas are exact, computationally testable, and closely connected to the repository's existing themes of Thue--Morse cancellation, $q$-binomial arithmetic, Bernoulli/Bell expansions, and Lambert-$W$ endpoint inversion.

The strongest immediate research opportunities are to formalize the complex-weight series and Newton interpolation, prove derivative-stable all-orders inverse asymptotics, develop the arithmetic of the rational Jacobi parameters and autocorrelation values, and use the Mellin--Laplace bridge to understand negative complex moments.

\appendix

\section{Selected exact formulas at a glance}

For convenience, the central formulas are collected here:
\begin{align*}
 I_0^nF(x)&=2^{\binom n2}F(x/2^n),\\
 \int_0^xt^pF(t)\dd t
 &=\sum_{j=0}^{p}(-1)^j\frac{p!}{(p-j)!}
 2^{\binom{j+1}{2}}x^{p-j}F(x/2^{j+1}),\\
 A_\nu(x)
 &=\sum_{j\ge0}(-1)^j\falling{\nu}{j}
 2^{\binom{j+1}{2}}x^{\nu-j}F(x/2^{j+1}),\\
 \mathcal M_F(z)&=\frac{1-D(z)}z,
 \qquad \mathcal M_{\upf,+}(z)=\frac{D(z)}z,\\
 D(z)&=1-z\sum_{j\ge0}(-1)^j\falling{z-1}{j}
 \frac{d_{j+1}}{(j+1)!},\\
 B(s)&=\prod_{j\ge1}\frac{1-e^{-s/2^j}}{s/2^j},\\
 \int_0^yG(t)^p\dd t&=yG(y)^p-pA_{p-1}(G(y)),\\
 G^{(n)}(y)&=\frac{Q_n[F'](G(y))}{F'(G(y))^{2n-1}},\\
 S(w)&=\int_0^1\frac{\dd y}{w-G(y)}
 =\sum_{n\ge0}\frac{d_n}{w^{n+1}}.
\end{align*}

\section{Bibliographic and repository notes}

The repository was inspected through its live GitHub directory pages and raw TeX views on 27 August 2026.  Because the draft and archive trees contain superseded copies, an archived duplicate was used to identify earlier claims but was not counted as an independent source.  The accompanying report is self-contained and does not require the repository to compile.

\begin{thebibliography}{99}

\bibitem{ProveItMain}
V. Reshetnikov,
\emph{ProveIt: Analysis/FabiusFunction documentation},
main monograph, Lean walkthrough, glossary, papers, archive, and draft frontier tree,
\url{https://github.com/VladimirReshetnikov/ProveIt/tree/main/Analysis/FabiusFunction/docs},
accessed 27 August 2026.

\bibitem{ProveItFrontiers}
V. Reshetnikov,
\emph{Non-formalized Research Frontiers for the Fabius and Rvachev Functions},
TeX synthesis in the \repo{} repository, version consulted 27 August 2026.

\bibitem{Fabius1966}
J. Fabius,
``A probabilistic example of a nowhere analytic $C^\infty$-function,''
\emph{Zeitschrift f\"ur Wahrscheinlichkeitstheorie und Verwandte Gebiete}
\textbf{5} (1966), 173--174.

\bibitem{Arias1982}
J. Arias de Reyna,
``An infinitely differentiable function with compact support: definition and properties,''
English translation of \emph{Rev. Real Acad. Ciencias Madrid} \textbf{76} (1982), 21--38,
\url{https://arxiv.org/abs/1702.05442}.

\bibitem{Arias2018}
J. Arias de Reyna,
``Arithmetic of the Fabius function,''
\emph{INTEGERS} \textbf{18} (2018), Paper A51,
\url{https://arxiv.org/abs/1702.06487}.

\bibitem{Haugland}
J. K. Haugland,
``Evaluating the Fabius function,''
\url{https://arxiv.org/abs/1609.07999}, latest version consulted August 2026.

\bibitem{Aistleitner}
C. Aistleitner, R. Hofer, and G. Larcher,
``On parametric Thue--Morse sequences and lacunary trigonometric products,''
\url{https://arxiv.org/abs/1502.06738}.

\bibitem{Toth}
L. T\'oth,
``Linear combinations of Dirichlet series associated with the Thue--Morse sequence,''
\url{https://arxiv.org/abs/2211.13570}.

\bibitem{BGT}
N. H. Bingham, C. M. Goldie, and J. L. Teugels,
\emph{Regular Variation},
Cambridge University Press, 1987; second printing with additions, 1989.

\bibitem{Comtet}
L. Comtet,
\emph{Advanced Combinatorics: The Art of Finite and Infinite Expansions},
D. Reidel, 1974.

\bibitem{Widder}
D. V. Widder,
\emph{The Laplace Transform},
Princeton University Press, 1941.

\bibitem{deBoor}
C. de Boor,
\emph{A Practical Guide to Splines}, revised edition,
Springer, 2001.

\bibitem{NIST}
NIST Digital Library of Mathematical Functions,
chapters on Bernoulli numbers, Bell polynomials, orthogonal polynomials, and asymptotic methods,
\url{https://dlmf.nist.gov/}.

\end{thebibliography}

\end{document}
